% 第 3 题解析：三种入射角下的光路（甲/乙/丙，TikZ）
% 甲：折射光直接由 B（圆弧中点）出射，折射角 45°；
% 乙：增大入射角，折射光射至 BC 弧，弧上入射角小于 45°（不可能全反射）；
% 丙：减小入射角，折射光射至 AB 弧，弧上入射角大于 45°（可能全反射）。
\begin{tikzpicture}[font=\small, scale=0.8]
  % ---------- 甲 ----------
  \begin{scope}[shift={(0,0)}]
    \coordinate (O1) at (0,0);
    \coordinate (A1) at (-2.4,0);
    \coordinate (C1) at (2.4,0);
    \coordinate (B1) at (0,-2.4);
    \draw[thick] (A1) -- (C1);
    \draw[thick] (A1) arc[start angle=180, end angle=360, radius=2.4];
    \draw[-{Stealth}] (-3.9,2.1) -- (A1);
    \draw[-{Stealth}] (A1) -- (B1);
    \node[below] at (0,-3.2) {甲};
  \end{scope}
  % ---------- 乙 ----------
  \begin{scope}[shift={(6.2,0)}]
    \coordinate (O2) at (0,0);
    \coordinate (A2) at (-2.4,0);
    \coordinate (C2) at (2.4,0);
    \draw[thick] (A2) -- (C2);
    \draw[thick] (A2) arc[start angle=180, end angle=360, radius=2.4];
    % 折射光线射至 BC 弧（更水平），入射角更大
    \coordinate (D2) at (1.15,-2.1);
    \draw[-{Stealth}] (-4.3,2.7) -- (A2);
    \draw[-{Stealth}] (A2) -- (D2);
    \node[below] at (0,-3.2) {乙};
  \end{scope}
  % ---------- 丙 ----------
  \begin{scope}[shift={(12.4,0)}]
    \coordinate (O3) at (0,0);
    \coordinate (A3) at (-2.4,0);
    \coordinate (C3) at (2.4,0);
    \draw[thick] (A3) -- (C3);
    \draw[thick] (A3) arc[start angle=180, end angle=360, radius=2.4];
    % 折射光线射至 AB 弧（更竖直），入射角更小
    \coordinate (E3) at (-1.15,-2.1);
    \draw[-{Stealth}] (-3.4,1.3) -- (A3);
    \draw[-{Stealth}] (A3) -- (E3);
    \node[below] at (0,-3.2) {丙};
  \end{scope}
\end{tikzpicture}
